In the HL-LHC era the expected instantaneous luminosities will be a factor of five larger than the current LHC nominal conditions. While these more intense conditions will enable the LHC experiments to accumulate datasets more than an order of magnitude larger than all of the preceding LHC runs, it will also exacerbate the experimental challenge of the multiple pile-up interactions appearing in each event. To mitigate degradation of LHC event reconstruction performance resulting from potentially hundreds of pile-up (PU) interactions expected in an HL-LHC event, both CMS and ATLAS will incorporate precision timing measurements into their detectors, significantly expanding the capabilities of the experiments through a combination of new sub-detectors and dedicated upgrades to existing detectors, such as the ECAL. This will facilitate the development of 4D reconstruction methods that will provide new tools for both pile-up suppression and a range of new experimental techniques to probe physics signatures and long-lived particles. 

The ECAL, through a combination of the scintillation timescale of PbWO$^{4}$ crystals, the electronic pulse shaping, and the sampling rate, allows for excellent time resolution to be obtained. The time reconstruction in ECAL is sensitive to the presence of pulses arising from particles produced in PU collisions. As the LHC moves to higher luminosity in the HL-LHC era, the effects of OOT-PU in the events become a significant issue in many areas.  The ECAL performance group has set a goal to achieve a $\sim$30 ps time resolution for HL-LHC data taking. As a significant step to prepare for operations in high OOT-PU environments, and to meet this goal, studies of methods to mitigate the effects of OOT-PU on the timing reconstruction have been performed.  Development of a ECAL timing reconstruction method that is OOT-PU aware will result in an improved time resolution that will not only improve the precision of many current physics measurements, it will also provide a foundation to meet the goal of the ECAL to achieve a $\sim$30 ps time resolution. The Cross  Correlation time reconstruction ( CC ) was developed to meet these challenges. The CC time reconstruction will not only support the incorporation of precision timing in the HL-LHC era, but will also allow us to begin to explore some of the possibilities of 4D reconstruction methods, that will facilitated by precision timing, in Run III.

The CC time reconstruction leverages the existing OOT pulse amplitudes determined by the multi-fit time reconstruction for the timing reconstruction. Given these OOT amplitudes and the known pulse shape of an isolated signal in a crystal, the nominal, in-time, pulse contribution in the measured amplitudes for a rechit is reconstructed. By extracting the nominal pulse to use in reconstruction of the rechit time and finding the time that, via a cross correlation function, best matches the extracted signal amplitudes to the pulse shape template, the effects of OOT-PU are mitigated, significantly improving the ECAL time resolution. 

In studies of average rechit time by BX, the CC time reconstruction significantly reduces dependence of the rechit time on the BX position, mitigating the impact of OOT-PU on the reconstruction, and greatly reduces the energy dependence in the time reconstruction. In direct comparison, the resolution of the times reconstructed by the CC time reconstruction in data show that the CC time reconstruction produces an improvement $\sim$30 ps for a same RO unit resolution and an improvement of $\sim$35 ps for the $Z\rightarrow e^{+}e^{-}$ resolution compared to the default ratio time reconstruction reconstruction in Run 2. In Run 3 the CC time reconstruction was implemented in CMSSW and an online time calibration for CC was deployed. A devoted full data reconstruction where CC was used as the default ECAL time reconstruction was produced for a selected data taking period and time resolution achieved was compared to the nominal prompt reconstruction for the data period. In these datasets it was seen that CC achieved an improvement of $\sim$80 ps for  the same RO unit resolution over the ratio time reconstruction. In year 2025 of Run 3 the CC time reconstruction was deployed as the default time reconstruction for data taking.  The CC time reconstruction has achieved SRU resolution of $\sim$135 ps and $Z\rightarrow e^{+}e^{-}$ resolution of $\sim$180 ps in prompt reconstruction in early 2025 data. With the utilization of the CC time reconstruction we expect comparable, if not improved, resolutions for all of Run III.

%The CC time reconstruction has achieved SRU resolution of $\sim$110 ps and $Z\rightarrow e^{+}e^{-}$ resolution of $\sim$140? ps in prompt reconstruction during initial deployment in Run 3 . With the utilization of the CC time reconstruction we expect comparable, if not improved, resolutions for all of Run III.

%The CC time reconstruction is expected to achieve $Z\rightarrow e^{+}e^{-}$, analysis relevant E/gamma object, resolutions of $\sim$100 ps, compared to $\sim$300 ps typical in Run II.  The CC time reconstruction has achieved SRU resolution of $\sim$110 ps and $Z\rightarrow e^{+}e^{-}$ resolution of $\sim$140? ps in prompt reconstruction during initial deployment in Run 3 . With the utilization of the CC time reconstruction we expect comparable, if not improved, resolutions for all of Run III.
